% Source: https://artofproblemsolving.com/community/c3h1479964_modular_arithmetic
% Stable identity: aops-topic-1479964-q01
\textbf{(a)} $10-(-1)=11$ is divisible by $11$, so $10\equiv-1\pmod{11}$. Also $5-10=-5$ and $10-0=10$ are both divisible by $5$, so both statements hold. A congruence does not say that one side is ``the remainder'' of the other: it says the two numbers leave the \emph{same} remainder, and $5$, $10$ and $0$ all leave remainder $0$ on division by $5$. Similarly $10$ and $-1$ both leave remainder $10$ on division by $11$ (since $-1=11\cdot(-1)+10$). So $5\equiv10\equiv0\pmod5$ are all true at once; congruence modulo $m$ is an equivalence relation, so these statements are consistent.

\textbf{(b)} Write $a-b=mk$ and $c-d=ml$ with integers $k,l$. Then
\[
(a+c)-(b+d)=m(k+l),
\]
so $a+c\equiv b+d$. For the product,
\[
ac-bd=ac-bc+bc-bd=c(a-b)+b(c-d)=m(ck+bl),
\]
so $ac\equiv bd\pmod m$. $\blacksquare$

Applying the product rule repeatedly, $a\equiv b$ implies $a^n\equiv b^n$ for every positive integer $n$. This is why congruences are so powerful: in any sum or product we may replace a number by anything congruent to it, for instance a small or negative number.

\textbf{(c)} From $10\equiv-1\pmod{11}$ we get $10^{2017}\equiv(-1)^{2017}=-1\equiv10\pmod{11}$. The remainder is $\boxed{10}$.
